\documentclass[11pt]{article}
\usepackage[margin=1in]{geometry}
\usepackage{amsmath,amssymb}
\usepackage{parskip}


\begin{document}

\begin{center}
{\Large\bfseries Macro HW2 --- Notes}\\[0.35em]
{\small ECON*6020 Macroeconomic Theory I, Assignment 2 (Arrow learning-by-investing)}
\end{center}

\hrule
\medskip

\textbf{Model.}
\begin{gather*}
\dot k = s\,(w + rk) - \delta k, \\
Y_i = (AL_i)^{1-\alpha}K_i^{\alpha}, \qquad
A = a\Big(\frac{1}{N}\sum_{i=1}^{N} K_i\Big)^{b}.
\end{gather*}
$L$ identical consumers each supply one unit of labour and rent capital to $N$ firms. Each firm has
constant returns in its own $K_i, L_i$, but technology $A$ rises with the \emph{average} capital
stock. Firms take $A$ as given, so capital accumulation creates an externality
(Arrow's learning-by-investing).

% ------------------------------------------------------------------
\section*{(a) Market clearing in the factor markets}

In each factor market, the total quantity demanded by firms equals the total quantity supplied by all consumers:
\[
L = \sum_{i=1}^{N} L_i, \qquad K = k\cdot L = \sum_{i=1}^{N} K_i .
\]
Labour: every consumer supplies one unit, so the firms together hire all $L$ workers.
Capital: the capital consumers own, $kL$, is all rented out to the firms.

% ------------------------------------------------------------------
\section*{(b) Real wage and rental rate}

Derive $r$ and $w$ as functions of the model parameters and the average capital--labour ratio.

Competitive firms pay each factor its marginal product, taking $A$ as given. Because all firms face
the same $w$, $r$ and $A$ and have constant returns, they all choose the same ratio
$K_i/L_i = K/L = k$. We then substitute $A = a(K/N)^b$ and $K = Lk$.

\textbf{Wage.}
\begin{align*}
w = \frac{\partial (AL_i)^{1-\alpha}K_i^{\alpha}}{\partial L_i}
  &= (1-\alpha)A^{1-\alpha}L_i^{-\alpha}K_i^{\alpha} \\
  &= (1-\alpha)\Big[a\Big(\frac{K}{N}\Big)^{b}\Big]^{1-\alpha}\Big(\frac{K_i}{L_i}\Big)^{\alpha}
   = (1-\alpha)a^{1-\alpha}\Big(\frac{K}{N}\Big)^{(1-\alpha)b}k^{\alpha} \\
  &= (1-\alpha)a^{1-\alpha}\Big(\frac{L\cdot k}{N}\Big)^{(1-\alpha)b}k^{\alpha} \\
  &= (1-\alpha)a^{1-\alpha}\Big(\frac{L}{N}\Big)^{(1-\alpha)b}k^{(1-\alpha)b+\alpha}.
\end{align*}

\textbf{Rental rate.} With $Y_i = (AL_i)^{1-\alpha}K_i^{\alpha}$ and $K = L\cdot k$:
\begin{align*}
r = \frac{\partial Y_i}{\partial K_i}
  &= \alpha A^{1-\alpha}L_i^{1-\alpha}K_i^{\alpha-1}
   = \alpha\Big\{a\Big[\Big(\frac{L\cdot k}{N}\Big)\Big]^{b}\Big\}^{1-\alpha}k^{\alpha-1} \\
  &= \alpha a^{1-\alpha}\Big(\frac{L}{N}\Big)^{b(1-\alpha)}k^{b(1-\alpha)+(\alpha-1)}
   = \alpha a^{1-\alpha}\Big(\frac{L}{N}\Big)^{b(1-\alpha)}k^{(1-\alpha)(b-1)}.
\end{align*}

% ------------------------------------------------------------------
\section*{(c) Increasing returns to scale for $b>0$}

Idea: add up output across firms to get aggregate output $Y$ as a function of $K$ and $L$, then
scale both inputs by $z$ and compare the result with $zY$.

\begin{align*}
Y_i &= (AL_i)^{1-\alpha}K_i^{\alpha}, \\
Y = \sum_{i=1}^{N}(AL_i)^{1-\alpha}K_i^{\alpha}
  &= \sum_{i=1}^{N}A^{1-\alpha}L_i\cdot L_i^{-\alpha}K_i^{\alpha}
   = \sum_{i=1}^{N}A^{1-\alpha}L_i\Big(\frac{K_i}{L_i}\Big)^{\alpha} \\
  &= \sum_{i=1}^{N}A^{1-\alpha}k^{\alpha}L_i
   = A^{1-\alpha}k^{\alpha}\sum_{i=1}^{N}L_i
   = A^{1-\alpha}k^{\alpha}L \\
  &= a^{1-\alpha}\bigg(\frac{\sum_{i=1}^{N}K_i}{N}\bigg)^{b(1-\alpha)}k^{\alpha}L
   = a^{1-\alpha}\Big(\frac{K}{N}\Big)^{b(1-\alpha)}k^{\alpha}L \\
  &= a^{1-\alpha}\Big(\frac{K}{N}\Big)^{b(1-\alpha)}L\Big(\frac{K}{L}\Big)^{\alpha}.
\end{align*}

If we multiply both inputs $K$ and $L$ by $z$, we get
\begin{align*}
&\phantom{=}\; a^{1-\alpha}\Big(\frac{zK}{N}\Big)^{b(1-\alpha)}zL\Big(\frac{zK}{zL}\Big)^{\alpha} \\
&= z^{\,b-b\alpha+1}\,a^{1-\alpha}\Big(\frac{K}{N}\Big)^{b(1-\alpha)}L\Big(\frac{K}{L}\Big)^{\alpha}
 = z^{\,1+b(1-\alpha)}\,Y.
\end{align*}
Since $0<\alpha<1$ (and $b>0$),
\[
0 < b\alpha < b \;\Longrightarrow\; b - b\alpha > 0 \;\Longrightarrow\; b - b\alpha + 1 > 1 .
\]
Therefore, for $z>1$, output rises by more than the factor $z$ by which $K$ and $L$ are increased:
the economy as a whole has increasing returns to scale in capital and labour.

Intuition: each firm alone has constant returns, but when all firms raise capital together the
average capital stock rises, which raises $A$ for everyone. That spillover gives the extra
$z^{b(1-\alpha)}$. If $b=0$ there is no spillover and we are back to constant returns.

% ------------------------------------------------------------------
\section*{(d) Transition dynamics to the steady state}

Plug $w$ and $r$ from (b) into the budget constraint $\dot k = s(w+rk)-\delta k$:
\begin{align*}
w + rk &= (1-\alpha)a^{1-\alpha}\Big(\frac{L}{N}\Big)^{(1-\alpha)b}k^{(1-\alpha)b+\alpha}
        + \alpha a^{1-\alpha}\Big(\frac{L}{N}\Big)^{b(1-\alpha)}k^{(1-\alpha)(b-1)+1} \\
       &= a^{1-\alpha}\Big(\frac{L}{N}\Big)^{b(1-\alpha)}
          \Big[(1-\alpha)k^{(1-\alpha)b+\alpha} + \alpha k^{(1-\alpha)b+\alpha}\Big] \\
       &= a^{1-\alpha}\Big(\frac{L}{N}\Big)^{b(1-\alpha)}k^{(1-\alpha)b+\alpha} = y .
\end{align*}
This matches output per worker from (c):
$y = Y/L = a^{1-\alpha}\big(\frac{kL}{N}\big)^{b(1-\alpha)}k^{\alpha}
     = a^{1-\alpha}\big(\frac{L}{N}\big)^{b(1-\alpha)}k^{b(1-\alpha)+\alpha}$.
So income per worker is all spent or saved, and
\[
\dot k = s y - \delta k = s a^{1-\alpha}\Big(\frac{L}{N}\Big)^{b(1-\alpha)}k^{b(1-\alpha)+\alpha} - \delta k .
\]
Dividing by $k$ gives the growth rate of capital per worker:
\[
\gamma_k \equiv \frac{\dot k}{k}
 = s a^{1-\alpha}\Big(\frac{L}{N}\Big)^{b(1-\alpha)}k^{(1-\alpha)(b-1)} - \delta .
\]
Everything depends on the sign of the exponent $(1-\alpha)(b-1)$. Since $y \propto k^{b(1-\alpha)+\alpha}$,
income per worker grows at $\gamma_y = [b(1-\alpha)+\alpha]\,\gamma_k$, so it follows the same pattern.

\textbf{(i) $b<1$: stable steady state.} The exponent is negative, so the growth rate falls as $k$
rises. Setting $\gamma_k=0$,
\[
k^* = \bigg[\frac{s a^{1-\alpha}(L/N)^{b(1-\alpha)}}{\delta}\bigg]^{\frac{1}{(1-\alpha)(1-b)}} .
\]
If $k<k^*$ then $\gamma_k>0$, and growth slows as $k$ approaches $k^*$; if $k>k^*$ then $k$ shrinks.
The economy converges to $k^*$ just as in the Solow model, and long-run growth is zero.

\textbf{(ii) $b>1$: unstable steady state.} The exponent is positive, so the growth rate
\emph{rises} with $k$. The same $k^*$ exists but is unstable. If $k>k^*$, growth is positive and keeps
speeding up (explosive growth). If $k<k^*$, $k$ shrinks faster and faster toward zero.
There is no convergence.

\textbf{(iii) $b=1$: no transition dynamics.} The exponent is zero, so $\gamma_k$ does not depend on
$k$. The economy grows at one constant rate from the start (an AK model); see (e).

% ------------------------------------------------------------------
\section*{(e) Steady-state growth rate for $b=1$}

With $b=1$, output per worker is linear in $k$:
$y = a^{1-\alpha}(L/N)^{1-\alpha}k$. So
\[
\gamma = \frac{\dot k}{k} = \frac{\dot y}{y} = s a^{1-\alpha}\Big(\frac{L}{N}\Big)^{1-\alpha} - \delta .
\]
This is the steady-state growth rate of capital and income per worker. With no population growth
it is also the growth rate of total output $Y$. Growth is positive if
$s a^{1-\alpha}(L/N)^{1-\alpha} > \delta$, and it rises with the savings rate $s$.
It also rises with $L/N$, the number of workers per firm (a scale effect).

\end{document}